\documentclass[10pt,a4paper]{amsart}
\usepackage[margin=19mm]{geometry}
\usepackage{amsmath,amssymb,mathtools}
\usepackage[scr=rsfs,cal=euler]{mathalfa}
\usepackage{xfrac}
\usepackage[unicode,hidelinks,bookmarks=false]{hyperref}
\newcommand{\ie}{i.e.\ }
\newcommand{\cf}{cf.\ }
\newcommand{\resp}{respectively\ }
\newcommand{\Subsection}{\subsection}
\providecommand{\nobreakdash}{\nobreak\hbox{-}}
\setlength{\emergencystretch}{2em}
\multlinegap0pt
\usepackage{amssymb}
\usepackage[shortlabels]{enumitem}
\usepackage{float}
\usepackage{mathtools}
\usepackage{xcolor}
\usepackage{tikz}
\newtheorem{lemma}{Lemma}
\title{Notes on Bilinear Bessel Potentials}
\author{Ana \v{C}olovi\'c, Xinyu Gao}
\date{Source: arXiv:2603.15962v1 (CC BY 4.0)}
\begin{document}
\maketitle

\noindent\emph{Source and licence.} Notes on Bilinear Bessel Potentials. Version arXiv:2603.15962v1 (CC BY 4.0), distributed by arXiv under the Creative Commons Attribution 4.0 International licence, http://creativecommons.org/licenses/by/4.0/, as recorded in the arXiv licence metadata for this exact version and shown on the abstract page https://arxiv.org/abs/2603.15962v1. Full licence notice: \texttt{licenses/arxiv-2603.15962-CC-BY-4.0.txt}.\par\smallskip

\noindent\textit{Editorial scope.} Counted unit: the article's lemma lem:interp-LPLQ (source0.tex lines 784--801). The article's complete original proof of it is delivered with it (source0.tex lines 803--909, one \texttt{proof} environment of 107 lines). No mathematics is written, repaired or re-derived: the statement and the proof are the article's own lines, in the article's own notation, and no retained line is substituted or re-flowed.  No further source-local context is needed: every cross-reference of every retained line resolves inside the two delivered spans.  Nothing was omitted from any retained span, so the package carries no omission record.  No retained line contains a citation, so the wrapper carries no bibliography and no bibliography entry.  No retained line contains a figure environment, an image-inclusion command or drawing code, so no image is delivered and the package declares no asset dependency.  Editorial formatting only, in addition: an AMS \texttt{amsart} preamble that restates the article's own declarations and macros used by the retained lines, namely the environments \texttt{lemma} on separate counters, together with the standard packages those lines need.  The retained environments print in this document as Lemma 1 in order of appearance, whereas the article prints the same items with the numbering of its own full document.  The complete native source of the article as distributed in the arXiv e-print for this exact version is bundled unchanged as \texttt{sources/196510/source0.tex}.
\begin{lemma}
\label{lem:interp-LPLQ}
Let $0<s<n$, $1\le p,q\le\infty$.
Then for every $0<\theta<1$ and every $0<\alpha\le\infty$,
\[
\|\mathcal J_s(f,g)\|_{L^{r_\theta,\alpha}}
\lesssim
A^{1-\theta}B^\theta\,\|f\|_{L^p}\|g\|_{L^q},
\]
where
\[
\frac1{r_\theta}
=
\frac{1-\theta}{r_1}
+
\frac{\theta}{r_2}.
\]
\end{lemma}

\begin{proof}
Let $h=\mathcal J_s(f,g)$.
By hypothesis,
\[
\|h\|_{L^{r_1}}
\le A\|f\|_{L^p}\|g\|_{L^q},
\qquad
\|h\|_{L^{r_2,\infty}}
\le B\|f\|_{L^p}\|g\|_{L^q}.
\]

Let $h^*$ denote the decreasing rearrangement of $h$.

From Chebyshev's inequality and the strong $L^{r_1}$ bound,
\[
h^*(t)
\le
\|h\|_{L^{r_1}}\,t^{-1/r_1}
\le
A\|f\|_{L^p}\|g\|_{L^q}\,t^{-1/r_1},
\qquad t>0.
\]
Similarly, from the weak-type bound,
\[
h^*(t)
\le
\|h\|_{L^{r_2,\infty}}\,t^{-1/r_2}
\le
B\|f\|_{L^p}\|g\|_{L^q}\,t^{-1/r_2},
\qquad t>0.
\]
Therefore,
\[
h^*(t)
\le
\min\Bigl(
A\|f\|_{L^p}\|g\|_{L^q}\,t^{-1/r_1},
\;
B\|f\|_{L^p}\|g\|_{L^q}\,t^{-1/r_2}
\Bigr).
\]

Let $0<\theta<1$ and define $r_\theta$ by
$\frac1{r_\theta}
=
\frac{1-\theta}{r_1}
+
\frac{\theta}{r_2}$.


Let $t_0>0$ be determined by
\[
A t_0^{-1/r_1}=B t_0^{-1/r_2},\qquad
t_0^{\,1/r_2-1/r_1}=\frac{A}{B}.
\]

Assume first that $0<\alpha<\infty$.
\begin{align*}
\|h\|_{L^{r_\theta,\alpha}}^\alpha&=
\int_0^\infty
\bigl(t^{1/r_\theta}h^*(t)\bigr)^\alpha
\frac{dt}{t}\\
&\le
\int_0^{t_0}
\bigl(B\|f\|_{L^p}\|g\|_{L^q}\,
t^{1/r_\theta-1/r_2}\bigr)^\alpha
\frac{dt}{t}\\
&+
\int_{t_0}^{\infty}
\bigl(A\|f\|_{L^p}\|g\|_{L^q}\,
t^{1/r_\theta-1/r_1}\bigr)^\alpha
\frac{dt}{t}.
\end{align*}

Since $r_1<r_2$
\[
\frac1{r_\theta}-\frac1{r_2}>0,
\qquad
\frac1{r_\theta}-\frac1{r_1}<0,
\]
both integrals converge.

A direct computation gives
\[
\|h\|_{L^{r_\theta,\alpha}}^\alpha
\le C
(A\|f\|_{L^p}\|g\|_{L^q})^{\alpha(1-\theta)}
(B\|f\|_{L^p}\|g\|_{L^q})^{\alpha\theta},
\]
where $C=C(n,s,p,q)$ a constant determined only on $n,s,p,q$.
Taking $\alpha$-th roots yields
\[
\|h\|_{L^{r_\theta,\alpha}}
\lesssim
A^{1-\theta}B^\theta
\|f\|_{L^p}\|g\|_{L^q}.
\]

When $\alpha=\infty$, we use
\[
\|h\|_{L^{r_\theta,\infty}}
=
\sup_{t>0} t^{1/r_\theta}h^*(t),
\]
together with the same pointwise bound.
The supremum is attained up to constants at $t\approx t_0$,
which gives the same estimate. \end{proof}

\end{document}
